% GitHub Repo and Documentation: https://github.com/celiobjunior/resume-template
% Copyright © 2025 Celio B Junior. All rights reserved.
% 
% Licensed under the Apache License, Version 2.0 (the "License");
% you may not use this file except in compliance with the License.
% You may obtain a copy of the License at
%
%     http://www.apache.org/licenses/LICENSE-2.0
%
% This template follows best practices from README.md

% Início do documento LaTeX: define tipo e formato do documento
% a4paper = tamanho A4, 10pt = tamanho base da fonte
\documentclass[a4paper,10pt]{article}

% --- PACOTES: BASE / IDIOMA ---
\usepackage[utf8]{inputenc}
\usepackage[T1]{fontenc}
\usepackage[portuguese]{babel}

% --- PACOTES: LAYOUT / TIPOGRAFIA ---
\usepackage{geometry}
\usepackage{parskip}
\usepackage{microtype}
\usepackage{enumitem}
\usepackage{titlesec}
\usepackage{array}
\usepackage{tabularx}

% --- PACOTES: LINKS / BOOKMARKS ---
\usepackage{hyperref}
\usepackage{bookmark}
\usepackage{xurl}

% --- CONFIGURAÇÃO: MARGENS / ESPAÇAMENTO ---
\geometry{top=1.5cm, bottom=1.5cm, left=1.5cm, right=1.5cm} % Margens da página
\setcounter{secnumdepth}{0}                                 % Remove numeração de seção
\setlist[itemize]{
    leftmargin=0.75em, % Recuo da lista (mais perto da margem da página)
    itemsep=0.2em,     % Espaço entre itens
    topsep=0.25em,     % Espaço antes/depois da lista
    parsep=0em,        % Espaço entre parágrafos no mesmo item
    partopsep=0em      % Espaço extra ao iniciar a lista
}

% Remove números de página e cabeçalhos para visual limpo do currículo
\pagestyle{empty}

% --- CONFIGURAÇÃO: METADADOS / LINKS ---
\hypersetup{
    pdftitle={CV Gustavo Silva Martins},   % Título do PDF
    pdfauthor={Gustavo Silva Martins},     % Autor do PDF
    colorlinks=true,          % Usa cor nos links (sem caixa)
    linkcolor=black,          % Cor de links internos
    urlcolor=black,           % Cor de URLs
    citecolor=black,          % Cor de citações
    bookmarksdepth=2          % Barra lateral: seção e subseção
}

% --- CONFIGURAÇÃO: TÍTULOS ---
\titleformat{\section}
{\Large\bfseries}
{}
{0em}
{}
[\titlerule\vspace{0.5ex}]
\titleformat{\subsection}
{\normalsize\bfseries}
{}
{0em}
{}
\titlespacing*{\subsection}{0pt}{0.35em}{0.15em}

% Macro para entradas do currículo com bookmark no PDF
\newcounter{cventry}
\newcommand{\cventry}[4]{%
    \refstepcounter{cventry}%
    \phantomsection%
    \pdfbookmark[2]{#1}{cventry-\thecventry}%
    \noindent\begin{tabularx}{\textwidth}{@{}>{\raggedright\arraybackslash}X >{\raggedleft\arraybackslash}X@{}}
    \textbf{#1} & #2 \\
    \textit{#3} & \textit{#4} \\
    \end{tabularx}
}

% --- INÍCIO DO DOCUMENTO ---
\begin{document}

% --- CABEÇALHO ---
% Substitua por suas informações pessoais
\begin{center}
    {\LARGE \textbf{Gustavo Silva Martins}}
    \\ [0.1cm]
    Capivari de Baixo, Santa Catarina
    {\textbullet}
    Email: \href{mailto:gustavosm994@gmail.com}{gustavosm994@gmail.com}
    {\textbullet}
    \href{https://linkedin.com/in/gustavo-martins994/}{linkedin.com/in/gustavo-martins994/}
    {\textbullet}
    \href{https://github.com/GusMartins499}{github.com/GusMartins499}
\end{center}

% --- SEÇÕES ---

\section{Atividades de Liderança}
    \cventry{FloripaBit}{Remoto}{Full-Stack Developer}{Jun-2023 - Atual}
        \begin{itemize}
            \item Apoio em code review, compreensão e aplicação das regras de negócio e auxílio na resolução de problemas para desenvolvedores juniores.
            \item Liderei a iniciativa de um queue-system para evitar gargalos de performance do atual sistema.
        \end{itemize}

% ------

\section{Experiência}
    \cventry{FloripaBit}{Remoto}{Full-Stack Developer}{Jun-2023 - Atual}
        \begin{itemize}
            \item Alcancei a redução significativa do tempo de geração de relatórios através da otimização de queries e o uso de streams resultando em uma melhor experiência para o usuário, reduzindo o tempo de geração do relatório em 50\%.
            \item Desenvolvi um queue-system usando workers do bullmq juntamente com o framework nestjs, resultando no não bloqueio do event loop do javascript.
            \item Colaborei com soluções para uma prova de conceito para uma plataforma de migração para o mercado de energia livre que resultou na primeira colocação do Hackathon da CCEE.
        \end{itemize}

    \cventry{Plathanus}{Palhoça - SC}{Full-Stack Developer}{Jan-2022 - Jun-2023}
        \begin{itemize}
            \item Alcancei o cargo de desenvolvedor pleno 1 após um período e de ter demonstrado progresso nas minhas habilidades técnicas, fui um case de sucesso no primeiro programa de inglês da empresa conquistando assim uma oportunidade para trabalhar em um projeto internacional (Projeto Tabas).
            \item Atuei na sustentação de projetos para grandes empresas como Professor Ferreto, Vialaser e no desenvolvimento de uma plataforma KPI para Sodexo com gráficos interativos e relatórios em xlsx e PDF.
        \end{itemize}

    \cventry{Cubo Sistemas}{Tubarão - SC}{Full-Stack Developer}{Jul-2020 - Set-2021}
        \begin{itemize}
            \item Conquistei o cargo de desenvolvedor júnior após um período trabalhando como suporte técnico e ter alcançado a senioridade máxima como suporte técnico e então promovido a desenvolvedor júnior onde atuei na construção de um sistema para emissão de nota fiscal eletrônica integrando com uma API open source em PHP.
        \end{itemize}

% ------

\section{Habilidades}
    \begin{itemize}
        \item \textbf{Técnicas:} Javascript/Typescript, Node.js, Nest.js, React, Next.js, Tailwindcss, Docker, Git, Cypress, Playwright, Svelte, A11y
        \item \textbf{Idiomas:} Português (Nativo), Inglês B2/C1
    \end{itemize}

% ------

% !!!!!!!!!!!!!!!!!!!!!!
% Mude esta seção para o começo, depois do cabeçalho e antes
% das suas experiências se você está procurando sua primeira
% vaga ou algum estágio. Do contrário, deixe como está.
% !!!!!!!!!!!!!!!!!!!!!!

\section{Educação}
    \cventry{UNISUL}{Tubarão - SC}{Bacharelado em Ciência da Computação}{2017 - julho de 2022}

\end{document}